\documentclass[11pt,a4paper]{article}
\usepackage[utf8]{inputenc}
\usepackage[T1]{fontenc}
\usepackage[french]{babel}
\usepackage{geometry}
\geometry{margin=2.5cm}
\usepackage{graphicx}
\usepackage{xcolor}
\usepackage{listings}
\usepackage{hyperref}
\usepackage{fancyhdr}
\usepackage{enumitem}
\usepackage{tcolorbox}
\usepackage{tikz}
\usepackage{amssymb}
\usepackage{booktabs}
\usepackage{array}
\usepackage{multirow}
\usepackage{pifont}
\newcommand{\faIcon}[1]{\ensuremath{\bullet}}

% Configuration des listes
\setlist[itemize]{leftmargin=*, itemsep=3pt, topsep=3pt}
\setlist[enumerate]{leftmargin=*, itemsep=3pt, topsep=3pt}

% Configuration du code
\lstset{
    basicstyle=\ttfamily\small,
    breaklines=true,
    frame=single,
    backgroundcolor=\color{gray!10},
    keywordstyle=\color{blue},
    commentstyle=\color{green!60!black},
    stringstyle=\color{orange},
    tabsize=4,
    showstringspaces=false
}

% En-tête et pied de page
\pagestyle{fancy}
\fancyhf{}
\fancyhead[L]{\small Architecture Dépôt Git – Phase 2}
\fancyhead[R]{\small Chatbot SUP'PTIC}
\fancyfoot[C]{\thepage}

% Couleurs personnalisées
\definecolor{suppticblue}{RGB}{0, 70, 140}
\definecolor{phase1color}{RGB}{180, 60, 60}
\definecolor{phase2color}{RGB}{30, 130, 80}
\definecolor{warningyellow}{RGB}{255, 193, 7}

% Environnements personnalisés
\newtcolorbox{important}[1][]{
  colback=red!5!white,
  colframe=red!75!black,
  fonttitle=\bfseries,
  title=#1
}

\newtcolorbox{info}[1][]{
  colback=blue!5!white,
  colframe=blue!75!black,
  fonttitle=\bfseries,
  title=#1
}

\newtcolorbox{astuce}[1][]{
  colback=green!5!white,
  colframe=green!75!black,
  fonttitle=\bfseries,
  title=#1
}

\newtcolorbox{phase1box}[1][]{
  colback=red!5!white,
  colframe=phase1color,
  fonttitle=\bfseries,
  title=#1
}

\newtcolorbox{phase2box}[1][]{
  colback=green!5!white,
  colframe=phase2color,
  fonttitle=\bfseries,
  title=#1
}

\newtcolorbox{decision}[1][]{
  colback=orange!5!white,
  colframe=orange!80!black,
  fonttitle=\bfseries,
  title=#1
}

\begin{document}

% ============================================================
% PAGE DE GARDE
% ============================================================
\begin{titlepage}
    \centering
\vspace*{1.5cm}

{\Huge \textbf{Architecture du Dépôt Git}} 

\vspace{0.4cm}
{\LARGE \textbf{Organisation des Branches -- Phase 2}}

\vspace{1.2cm}

{\Large Projet Chatbot SUP'PTIC}

\vspace{1cm}

\begin{minipage}{0.85\textwidth}
    \centering
    \normalsize
    Analyse de la situation héritée de la Phase 1, \\
    recommandations sur les branches existantes \\
    et nouvelle organisation pour la Phase 2
\end{minipage}



\end{titlepage}

\tableofcontents
\clearpage

% ─── Section 1 : Bilan Phase 1 ────────────────────────────────────────────────
\section{Bilan du Dépôt — Héritage de la Phase 1}
\label{sec:bilan}

\subsection{Rappel : Branches créées en Phase 1}

À l'issue de la Phase 1 (8-13 février 2026), le dépôt GitHub contenait les branches suivantes :

\vspace{0.5cm}
\begin{center}
\begin{tabular}{>{\ttfamily}l l l}
\toprule
\textbf{Branche} & \textbf{Type} & \textbf{Rôle en Phase 1} \\
\midrule
main     & Protégée & Version stable finale, prête pour démo \\
dev      & Protégée & Intégration de toutes les features \\
backend  & Équipe   & API Django REST + algorithme TF-IDF \\
frontend & Équipe   & Interface HTML/CSS/JS \\
database & Équipe   & Modèles Django, migrations, scripts \\
data     & Équipe   & Génération et nettoyage CSV (1000+ Q/R) \\
\bottomrule
\end{tabular}
\end{center}

\vspace{0.5cm}

\begin{phase1box}[Bilan Phase 1 — Ce qui a été livré]
\begin{itemize}
    \item[\ding{52}] Recherche TF-IDF avec similarité cosinus opérationnelle
    \item[\ding{52}] Base de données Django avec 1000+ FAQ chargées
    \item[\ding{52}] API REST complète (5 endpoints principaux)
    \item[\ding{52}] Interface web responsive + système de feedback
    \item[\ding{52}] Application PWA avec fonctionnalités offline
    \item[\ding{52}] Documentation complète (backend, frontend, PWA)
\end{itemize}
\end{phase1box}

\subsection{Répartition des équipes — Phase 1 vs Phase 2}

\vspace{0.3cm}
\begin{center}
\begin{tabular}{l c c l}
\toprule
\textbf{Équipe} & \textbf{Phase 1} & \textbf{Phase 2} & \textbf{Évolution} \\
\midrule
Backend           & 2 membres & 3 membres & \textcolor{phase2color}{$\uparrow$ +1} \\
Base de Données   & 2 membres & 1 membre  & \textcolor{phase1color}{$\downarrow$ -1} \\
Structuration Données & 4 membres & 4 membres & \textcolor{gray}{= Stable} \\
Frontend          & 2 membres & 2 membres & \textcolor{gray}{= Stable} \\
\midrule
\textbf{Total}    & \textbf{10} & \textbf{10} & — \\
\bottomrule
\end{tabular}
\end{center}

\vspace{0.5cm}

\begin{info}[Analyse de la nouvelle répartition]
Le changement principal entre les deux phases est le \textbf{renforcement du Backend} (+1 membre) compensé par une \textbf{réduction de l'équipe Base de Données} (-1 membre). Cette décision est cohérente avec l'état du projet : les modèles de données étant désormais stables, les efforts doivent se concentrer sur les nouvelles fonctionnalités applicatives.
\end{info}

\clearpage

% ─── Section 2 : Décision sur les branches Phase 1 ───────────────────────────
\section{Que Faire des Branches de la Phase 1 ?}
\label{sec:decision}

\subsection{Analyse branche par branche}

Avant de définir l'architecture de la Phase 2, il convient d'examiner chaque branche héritée et de décider de son sort.

\vspace{0.5cm}

\subsubsection{Branches protégées : \texttt{main} et \texttt{dev}}

\begin{astuce}[\ding{110} À conserver impérativement]
\begin{itemize}
    \item \textbf{\texttt{main}} — Contient la version stable 1.0.0. À garder comme référence historique et point de départ de la Phase 2.
    \item \textbf{\texttt{dev}} — Sert de branche d'intégration continue. Elle est réutilisée telle quelle en Phase 2.
\end{itemize}
Ces deux branches ne doivent \textbf{jamais} être supprimées.
\end{astuce}

\subsubsection{Branches d'équipe : \texttt{backend}, \texttt{frontend}, \texttt{database}, \texttt{data}}

\begin{decision}[\ding{110} Décision recommandée pour les branches d'équipe]
\textbf{Proposition : Archiver, puis recréer proprement.}

\vspace{0.3cm}
Les branches d'équipe de la Phase 1 portent l'historique de commits de tout le sprint précédent. Continuer à travailler directement dessus en Phase 2 risque de :
\begin{itemize}
    \item[\ding{56}] Créer de la confusion entre les travaux des deux phases
    \item[\ding{56}] Polluer l'historique Git avec des commits mélangés
    \item[\ding{56}] Rendre difficile la revue de code phase par phase
\end{itemize}

\vspace{0.3cm}
\textbf{Procédure recommandée :}
\begin{enumerate}
    \item Créer un \textbf{tag d'archivage} sur chaque branche Phase 1
    \item Supprimer les anciennes branches sur GitHub
    \item Recréer des branches fraîches pour la Phase 2 (à partir de \texttt{dev})
\end{enumerate}
\end{decision}

\subsubsection{Procédure d'archivage des branches Phase 1}

\begin{lstlisting}[language=bash, caption=Archivage des branches Phase 1 avant suppression]
# Etape 1 : Se positionner sur chaque branche et creer un tag
git checkout backend
git tag archive/phase1/backend
git push origin archive/phase1/backend

git checkout frontend
git tag archive/phase1/frontend
git push origin archive/phase1/frontend

git checkout database
git tag archive/phase1/database
git push origin archive/phase1/database

git checkout data
git tag archive/phase1/data
git push origin archive/phase1/data

# Etape 2 : Supprimer les branches distantes
git push origin --delete backend
git push origin --delete frontend
git push origin --delete database
git push origin --delete data

# Etape 3 : Nettoyer les branches locales
git branch -D backend
git branch -D frontend
git branch -D database
git branch -D data
\end{lstlisting}

\clearpage

% ─── Section 3 : Architecture Phase 2 ────────────────────────────────────────
\section{Architecture des Branches — Phase 2}
\label{sec:architecture}

\subsection{Vue d'ensemble}

\begin{phase2box}[Nouvelle architecture — Phase 2]
\begin{itemize}
    \item \textbf{2 branches permanentes protégées} : \texttt{main}, \texttt{dev}
    \item \textbf{4 branches d'équipe} recréées proprement depuis \texttt{dev}
    \item Même workflow qu'en Phase 1 (Pull Requests vers \texttt{dev})
\end{itemize}
\end{phase2box}

\vspace{0.5cm}

\begin{figure}[h]
\centering
\begin{tikzpicture}[node distance=1.8cm, auto]

    % Styles
    \tikzstyle{protected} = [rectangle, rounded corners, minimum width=3.2cm,
        minimum height=0.9cm, text centered, draw=suppticblue,
        fill=blue!15, font=\bfseries\small]
    \tikzstyle{team} = [rectangle, rounded corners, minimum width=3.2cm,
        minimum height=0.9cm, text centered, draw=phase2color,
        fill=green!10, font=\small]
    \tikzstyle{arrow} = [thick,->,>=stealth,color=gray!70]
    \tikzstyle{prArrow} = [thick,dashed,->,>=stealth,color=phase2color]

    % Branches protégées
    \node (main) [protected] at (12, 0) {main};
    \node (dev)  [protected] at (4, 0) {dev};

    % Branches équipes (en dessous)
    \node (backend)  [team] at (0,  -2.5) {backend};
    \node (frontend) [team] at (3.5,  -2.5) {frontend};
    \node (database) [team] at (7,  -2.5) {database};
    \node (data)     [team] at (10.5,-2.5) {data};

    % Flèches création depuis dev
    \draw [arrow] (dev) -- (backend);
    \draw [arrow] (dev) -- (frontend);
    \draw [arrow] (dev) -- (database);
    \draw [arrow] (dev) -- (data);

    % Pull Requests vers dev
    \draw [prArrow] (backend)  to[bend left=20] node[left,font=\tiny] {PR} (dev);
    \draw [prArrow] (frontend) to[bend left=10] node[right,font=\tiny] {PR} (dev);
    \draw [prArrow] (database) to[bend right=10] node[right,font=\tiny] {PR} (dev);
    \draw [prArrow] (data)     to[bend right=20] node[right,font=\tiny] {PR} (dev);

    % dev vers main
    \draw [arrow, color=suppticblue, very thick]
        (dev) -- node[above,font=\small\bfseries] {Validation finale} (main);

    % Légende
    \node[font=\footnotesize, draw=suppticblue, fill=blue!15, rounded corners,
          minimum width=2cm, minimum height=0.5cm] at (8, -4.5)
          {\textcolor{suppticblue}{\textbf{Protégée}}};
    \node[font=\footnotesize, draw=phase2color, fill=green!10, rounded corners,
          minimum width=2cm, minimum height=0.5cm] at (8, -5.5)
          {\textcolor{phase2color}{\textbf{Branche équipe}}};

\end{tikzpicture}
\caption{Flux de travail Git — Phase 2}
\end{figure}

\subsection{Création des branches Phase 2}

\begin{lstlisting}[language=bash, caption=Initialisation des branches Phase 2 depuis dev]
# S'assurer d'etre a jour
git checkout dev
git pull origin dev

# Creer les nouvelles branches Phase 2 depuis dev
git checkout -b backend
git push origin backend

git checkout dev
git checkout -b frontend
git push origin frontend

git checkout dev
git checkout -b database
git push origin database

git checkout dev
git checkout -b data
git push origin data
\end{lstlisting}

\clearpage

% ─── Section 4 : Arborescence complète du dépôt — Phase 2 ─────────────────────
\section{Arborescence complète du dépôt — Phase 2}
\label{sec:arborescence}

\begin{info}[Légende]
Les éléments marqués \texttt{\# [NEW]} sont nouveaux par rapport à la Phase 1. \\
Les dossiers \texttt{backend/rag\_data/} contiennent des artefacts générés localement et ne sont pas commités (dans \texttt{.gitignore}).
\end{info}

\vspace{0.3cm}
\begin{lstlisting}[basicstyle=\ttfamily\small, caption=Arborescence Phase 2, breaklines=true]
ChatBot/                          # Racine
|
+-- .venv/                        # Env virtuel Python
+-- .gitignore / requirements.txt # Dépendances (Phase 2 ajoutées)
+-- README.md
|
+-- backend/                      # Django
|   +-- config/                   # Paramètres + routes [NEW]
|   +-- faq/                      # FAQ (hérité)
|   +-- chatbot/                  # Moteur RAG
|   |   +-- views.py              # /api/chatbot/ask/ [NEW]
|   |   +-- engine/ [NEW]         # Pipeline RAG
|   |       +-- embedder.py       # MiniLM
|   |       +-- faiss_search.py   # FAISS
|   |       +-- llm_client.py     # Ollama Phi-3
|   |       +-- rag_pipeline.py   # Orchestration
|   |       +-- tfidf_fallback.py # Bascules TF-IDF
|   |       \-- prompt_builder.py # Prompt LLM
|   +-- users/                    # Utilisateurs + feedback
|   +-- rag_data/ [NEW]           # Artefacts FAISS
|   +-- scripts/ [NEW]            # build_index, reload_index, test_rag
|   +-- doc/                      # Docs backend
|   +-- manage.py / db.sqlite3
|
+-- data/                         # Données FAQ JSON
|   +-- faq_*.json                # Admissions, frais, etc.
|   +-- scripts/                  # validate_json, stats_dataset [NEW]
|   \-- doc/
|
+-- frontend/                     # Interface
|   +-- react-app/ [NEW]          # App React
|   +-- pwa/                      # Ancienne PWA (archivée)
|   +-- README.md / start_server.py
|
+-- docs/                         # Documentation générale
|   +-- architecture/
|   \-- setup/                    # OLLAMA_SETUP.md [NEW]
                             
\end{lstlisting}

%\clearpage

% ─── Conclusion ───────────────────────────────────────────────────────────────

\vfill

\begin{center}
\textit{Document préparé par :} \\
\textbf{NKOUMOU TJADE} \\
Chef de la Division de Programmation \\
Club Informatique SUP'PTIC

\end{center}

\end{document}